%! Rational Functions and Vertical Asymptotes — Unit 1, Lesson 1.9 — Group Activity
\documentclass[10pt]{article}
\usepackage{precalculus-article}
\usepackage{precalculus-boxes}
\newcommand{\TallMath}[1]{$\displaystyle #1\rule[-1.4em]{0pt}{3.2em}$}

\begin{document}

\pageheader{Unit 1, Lesson 1.9}{Group Activity}
\namepartnerperiod

\vspace{0.1in}

\begin{scenariobox}[Analyzing Pressure Gauges]{sky}
\small
Engineers monitor pressure in a hydraulic system. The pressure at different points in the
system is modeled by rational functions. A pressure ``spike'' --- where pressure grows without
bound --- corresponds to a vertical asymptote of the model. For each function below, determine
where pressure spikes occur.

\medskip
\renewcommand{\arraystretch}{1.4}
\begin{tabular}{@{} l l @{\qquad} l l @{}}
$r_1(x) = \dfrac{3}{x - 5}$ &
$r_2(x) = \dfrac{x + 2}{x^2 - 4}$ &
$r_3(x) = \dfrac{x^2 - 1}{(x-1)(x+3)}$ &
$r_4(x) = \dfrac{x}{(x+2)^3}$
\end{tabular}
\end{scenariobox}

\vspace{0.1in}

\begin{tcolorbox}[colback=white, colframe=black!40,
    title=\textbf{Tier R --- Remediate},
    fonttitle=\bfseries, arc=2mm, left=3mm, right=3mm, top=2mm, bottom=2mm]

For each function, set the denominator equal to zero and solve. Assume the numerator is not
zero at those points and write the vertical asymptote equation(s).

\begin{enumerate}[itemsep=10pt]
  \item $r_1(x) = \dfrac{3}{x - 5}$

        Denominator $= 0$ when $x =$ \blank{2cm}. \quad Vertical asymptote: \blank{3cm}

  \item $r_2(x) = \dfrac{x + 2}{x^2 - 4}$

        Factor the denominator: $x^2 - 4 =$ \blank{4cm}

        Denominator $= 0$ when $x =$ \blank{3.5cm}. \quad Vertical asymptote(s): \blank{3.5cm}

  \item $r_3(x) = \dfrac{x^2 - 1}{(x-1)(x+3)}$

        Denominator $= 0$ when $x =$ \blank{3.5cm}. \quad Vertical asymptote(s): \blank{3.5cm}

  \item $r_4(x) = \dfrac{x}{(x+2)^3}$

        Denominator $= 0$ when $x =$ \blank{2cm}. \quad Vertical asymptote: \blank{3cm}
\end{enumerate}
\end{tcolorbox}

\vspace{0.1in}

\begin{tcolorbox}[colback=white, colframe=black!40,
    title=\textbf{Tier A --- Apply},
    fonttitle=\bfseries, arc=2mm, left=3mm, right=3mm, top=2mm, bottom=2mm]

Factor both the numerator and denominator of each function. Determine whether each denominator
zero gives a \textbf{vertical asymptote} (VA) or a \textbf{potential hole} (if the numerator
is also zero at that point). Write the correct VAs.

\begin{enumerate}[itemsep=12pt]
  \item $r_1(x) = \dfrac{3}{x - 5}$: \quad
        Numerator at $x=5$: \blank{2cm}. \quad Result: \blank{4cm}

  \item $r_2(x) = \dfrac{x + 2}{x^2 - 4} = \dfrac{x+2}{(x-2)(x+2)}$

        At $x = 2$: numerator $=$ \blank{1.5cm} \quad $\Rightarrow$ \blank{4cm}

        At $x = -2$: numerator $=$ \blank{1.5cm} \quad $\Rightarrow$ \blank{4cm}

  \item $r_3(x) = \dfrac{x^2-1}{(x-1)(x+3)}$.
        Factor the numerator: $x^2 - 1 =$ \blank{4cm}

        At $x = 1$: numerator $=$ \blank{1.5cm} \quad $\Rightarrow$ \blank{4cm}

        At $x = -3$: numerator $=$ \blank{1.5cm} \quad $\Rightarrow$ \blank{4cm}

  \item $r_4(x) = \dfrac{x}{(x+2)^3}$: \quad
        At $x = -2$: numerator $=$ \blank{1.5cm} \quad $\Rightarrow$ \blank{4cm}
\end{enumerate}
\end{tcolorbox}

\vspace{0.1in}

\begin{tcolorbox}[colback=white, colframe=black!40,
    title=\textbf{Tier E --- Extend},
    fonttitle=\bfseries, arc=2mm, left=3mm, right=3mm, top=2mm, bottom=2mm]

\begin{enumerate}[itemsep=12pt]
  \item \textbf{Explaining a hole.} For $r_3(x) = \dfrac{x^2 - 1}{(x-1)(x+3)}$,
        factor the numerator as $(x-1)(x+1)$.

        \begin{enumerate}[label=(\alph*), itemsep=5pt]
          \item Why is $x = 1$ \emph{not} a vertical asymptote? Use the words
                ``multiplicity'' or ``cancels'' in your explanation.

                \vspace{0.25in}
                \writeline\\[4pt]
                \writeline

          \item Confirm that $x = -3$ \textbf{is} a vertical asymptote. Show your reasoning.

                \vspace{0.25in}
                \writeline
        \end{enumerate}

  \item \textbf{One-sided limits.} For $r_4(x) = \dfrac{x}{(x+2)^3}$, analyze the sign
        of $r_4(x)$ as $x \to -2$ from each side.

        \begin{enumerate}[label=(\alph*), itemsep=6pt]
          \item When $x$ is slightly greater than $-2$ (e.g., $x = -1.9$):
                Numerator sign: \blank{2cm}. \quad $(x+2)^3$ sign: \blank{2cm}.
                \quad So $r_4(x)$ is: \blank{2cm}

          \item When $x$ is slightly less than $-2$ (e.g., $x = -2.1$):
                Numerator sign: \blank{2cm}. \quad $(x+2)^3$ sign: \blank{2cm}.
                \quad So $r_4(x)$ is: \blank{2cm}

          \item Write the one-sided limits:
                \[
                  \lim_{x \to -2^+} r_4(x) = \blank{2.5cm}
                  \qquad\qquad
                  \lim_{x \to -2^-} r_4(x) = \blank{2.5cm}
                \]
        \end{enumerate}
\end{enumerate}
\end{tcolorbox}

\end{document}
